

The natural generalization of the polytopes in \Cref{eg:polytope:toric} is the moment polytopes of line bundles, which depends on the choice of linearizations.

\begin{definition}[c.f., {\cite[\S2]{KKLV_loc_prop_alggrp_act}}]\label[definition]{def:linz}
	Let $X$ be a $G$-variety, and $L$ a line bundle on $X$ with projection $\pi : L \to X$. A \emph{$G$-linearization} $\alpha$ of $L$ is an action of $G$ on $L$ (as a variety), such that $\pi$ is $G$-equivariant and $\alpha$ is linear on each fiber.

	A $G$-linearization $\alpha$ on the line bundle $L$ defines a $G$-linearization on $m L$ for every $m \in \mathbb{N}$. More generally, if $L$ is a $\mathbb{Q}$-line bundle, a \emph{$\mathbb{Q}$-$G$-linearization} $\alpha$ is a $G$-linearization of $m L$ for some $m \in \mathbb{N}(L)$; two $\mathbb{Q}$-$G$-linearizations $\alpha, \alpha'$ of $L$ are the same if they agree on some common multiple $m L$.

	Let $\mathbb{N}(L, \alpha) := \{m \in \mathbb{N}(L) \mid \alpha \text{ is a } G\text{-linearization of the line bundle } m L\}$; then, for every $m \in \mathbb{N}(L, \alpha)$, the vector space $H^0(X, m L)$ is a non-zero $G$-module.
\end{definition}
\begin{remark*}
	By \cite[\S2.4]{KKLV_loc_prop_alggrp_act}, for every line bundle $L$, there is a $G$-linearization for some multiple $m L$ ($X$ need to be normal); equivalently, every $\mathbb{Q}$-line bundle admits a $\mathbb{Q}$-$G$-linearization.
\end{remark*}

From now on, we will not distinguish the $G$-linearizations and the $\mathbb{Q}$-$G$-linearizations.

\begin{lemdefn}[c.f., {\cite[\S2.1]{CFPvS_mome_poly}}, based on \cite{Brion_sur_mome}]\label[lemdefn]{def:mome_poly}
	Given a polarized $G$-variety $(X, L)$ and a $G$-linearization $\alpha$, for every $m \in \mathbb{N}(L, \alpha)$, let
	\begin{equation*}
		P_m(X, L, \alpha) = \{\lambda \in \mathfrak{X}(B) \mid H^0(X, m L)_\lambda^{(B)} \neq \{0\} \}
	\end{equation*}
	(see \Cref{sec:prelim:sph} for the definition of $V_\lambda^{(B)}$), the \emph{moment polytope},
	\begin{equation*}
		P(X, L, \alpha) = \bigcup_{m \in \mathbb{N}(L, \alpha)} {1 \over m} P_m(X, L, \alpha) \subseteq \mathfrak{X}(B) \otimes \mathbb{Q}
	\end{equation*}
	is a convex hull of finitely many points.
\end{lemdefn}

The following result is directly deduced from the definition of moment polytope, 

\begin{lemma}\label[lemma]{lem:moneSPAN}
	For every $\lambda \in P(X, L, \alpha)$, the translation subset $P(X, L, \alpha) - \lambda \subseteq M_\mathbb{Q}$ spans $M_\mathbb{Q}$ (as a $\mathbb{Q}$-vector space). More specifically, when $m \in \mathbb{N}(L, \alpha)$ and line bundle $m L$ is very ample, for every $\lambda \in P_m(X, L, \alpha)$, the subset $P_m(X, L, \alpha) - \lambda$ generates $M$ (as an Abelian group).
\end{lemma}
\begin{proof}
	The first part is a reformulation of \cite[\S1.2, Proposition 3]{Brion_var_sph} (see also \cite[Lemma 2.21 (ii)]{Brion_intro_act_alg_grp}). We now prove the second part.

	For every $m \in \mathbb{N}(L, \alpha)$, we pick a nonzero section
	$s \in H^0(X, mL)^{(B)}_\lambda$, where the $\lambda \in P_m(X, L, \alpha)$. Now, for every
	$\lambda' \in P_m(X, L, \alpha)$, there exists a nonzero section
	$s' \in H^0(X, mL)^{(B)}_{\lambda'}$.
	The ratio of these two sections is a rational function, namely
	$f := s' / s \in \mathbb{C}(X)$, and it satisfies
	\begin{equation*}
		b \cdot f = (b \cdot s') / (b \cdot s)
		= (\lambda' - \lambda)(b)\, f,
	\end{equation*}
	for every $b \in B$. Therefore, by the definition of $M$, we have
	$\lambda' - \lambda \in M$, and thus
	$P_m(X, L, \alpha) - \lambda \subseteq M$.

	Now assume that $mL$ is very ample. This implies that for every
	$\Gamma \in P(X, L, \alpha)$, there exists a sufficiently large integer
	$N$, which is a multiple of $m$, such that
	\begin{equation*}
		N \Gamma = \sum_{{\lambda'} \in P_m(X, L, \alpha)} a_{\lambda'} {\lambda'},
	\end{equation*}
	where $a_{\lambda'} \in \mathbb{N}$ for every
	${\lambda'} \in P_m(X, L, \alpha)$.
	Finally, for every $u \in M$, we can find a sufficiently large $N$
	and $\Gamma, \Gamma' \in P(X, L, \alpha)$ such that
	$u = N(\Gamma - \Gamma')$, which can be further written as a
	$\mathbb{Z}$-linear combination of elements of
	$P_m(X, L, \alpha) - \lambda$.
\end{proof}